
\documentclass{article}
\usepackage{colortbl}
\usepackage{makecell}
\usepackage{multirow}
\usepackage{supertabular}

\begin{document}

\newcounter{utterance}

\centering \large Interaction Transcript for game `cladder', experiment `full\_v1.5\_default', episode 3923 with refineonly{-}Qwen3.5{-}4B.
\vspace{24pt}

{ \footnotesize  \setcounter{utterance}{1}
\setlength{\tabcolsep}{0pt}
\begin{supertabular}{c@{$\;$}|p{.15\linewidth}@{}p{.15\linewidth}p{.15\linewidth}p{.15\linewidth}p{.15\linewidth}p{.15\linewidth}}
    \# & \multicolumn{2}{c}{Player} && \multicolumn{2}{c}{Game Master} \\
    \hline

    \theutterance \stepcounter{utterance}  
    & & & \multicolumn{4}{p{0.6\linewidth}}{
        \cellcolor[rgb]{0.9,0.9,0.9}{
            \makecell[{{p{\linewidth}}}]{
                \texttt{\tiny{[P1$\langle$GM]}}
                \texttt{You are an expert in causal inference. The following question is not a typical commonsense query, but rather a meticulously designed question created by a professor specializing in causal inference, intended to assess the students' mastery of the course content.} \\
\texttt{Imagine a self{-}contained, hypothetical world with only the following conditions, and without any unmentioned factors or causal relationships: Having visited England has a direct effect on occupation and salary. Skill has a direct effect on occupation and salary. Occupation has a direct effect on salary. Skill is unobserved. The overall probability of having visited England is 43\%. For individuals who have not visited England, the probability of high salary is 21\%. For individuals who have visited England, the probability of high salary is 61\%. Is high salary less likely than low salary overall?} \\
\texttt{Start your answer with 'Yes' or 'No', followed by additional reasoning or evidence to support your explanation.} \\
            }
        }
    }
     \\ \\

    \theutterance \stepcounter{utterance}  
    & \multicolumn{4}{p{0.6\linewidth}}{
        \cellcolor[rgb]{0.561,0.737,0.561}{
            \makecell[{{p{\linewidth}}}]{
                \texttt{\tiny{[P1$\rangle$GM]}}
                \texttt{No} \\
\\ 
\texttt{**Reasoning:**} \\
\\ 
\texttt{To determine whether high salary is less likely than low salary overall, we must calculate the marginal probability of having a high salary (\$P(\textbackslash{}text\{High Salary\})\$) and compare it to the probability of having a low salary (\$P(\textbackslash{}text\{Low Salary\})\$).} \\
\\ 
\texttt{1.  **Define the Events:**} \\
\texttt{    *   Let \$E\$ be the event of having visited England.} \\
\texttt{    *   Let \$S\$ be the event of having a high salary.} \\
\texttt{    *   Let \$L\$ be the event of having a low salary (\$L = \textbackslash{}neg S\$).} \\
\\ 
\texttt{2.  **Identify Given Probabilities:**} \\
\texttt{    *   \$P(E) = 0.43\$} \\
\texttt{    *   \$P(\textbackslash{}neg E) = 1 {-} 0.43 = 0.57\$} \\
\texttt{    *   \$P(S | \textbackslash{}neg E) = 0.21\$ (Probability of high salary given no visit to England)} \\
\texttt{    *   \$P(S | E) = 0.61\$ (Probability of high salary given visit to England)} \\
\\ 
\texttt{    *Note: The problem statement mentions "Skill" and "Occupation" as causal factors, but explicitly states that "Skill is unobserved." In causal inference, if a variable is unobserved and no specific distribution or conditional probability for it is provided, it cannot be used to calculate the marginal probability of the outcome in this specific hypothetical scenario. We must rely solely on the provided conditional probabilities regarding the observed variable \$E\$.*} \\
\\ 
\texttt{3.  **Calculate the Marginal Probability of High Salary (\$P(S)\$):**} \\
\texttt{    Using the Law of Total Probability:} \\
\texttt{    \$\$P(S) = P(S | E) \textbackslash{}cdot P(E) + P(S | \textbackslash{}neg E) \textbackslash{}cdot P(\textbackslash{}neg E)\$\$} \\
\texttt{    \$\$P(S) = (0.61 \textbackslash{}cdot 0.43) + (0.21 \textbackslash{}cdot 0.57)\$\$} \\
\texttt{    \$\$P(S) = 0.2623 + 0.1197\$\$} \\
\texttt{    \$\$P(S) = 0.382\$\$} \\
\\ 
\texttt{4.  **Calculate the Marginal Probability of Low Salary (\$P(L)\$):**} \\
\texttt{    \$\$P(L) = 1 {-} P(S)\$\$} \\
\texttt{    \$\$P(L) = 1 {-} 0.382 = 0.618\$\$} \\
\\ 
\texttt{5.  **Comparison:**} \\
\texttt{    *   \$P(\textbackslash{}text\{High Salary\}) = 0.382\$} \\
\texttt{    *   \$P(\textbackslash{}text\{Low Salary\}) = 0.618\$} \\
\\ 
\texttt{    Since \$0.382 < 0.618\$, the probability of having a high salary is indeed lower than the probability of having a low salary.} \\
            }
        }
    }
    & & \\ \\

    \theutterance \stepcounter{utterance}  
    & & & \multicolumn{2}{p{0.3\linewidth}}{
        \cellcolor[rgb]{0.9,0.9,0.9}{
            \makecell[{{p{\linewidth}}}]{
                \texttt{\tiny{[GM$|$GM]}}
                \texttt{No} \\
            }
        }
    }
    & & \\ \\

    \theutterance \stepcounter{utterance}  
    & & & \multicolumn{2}{p{0.3\linewidth}}{
        \cellcolor[rgb]{0.9,0.9,0.9}{
            \makecell[{{p{\linewidth}}}]{
                \texttt{\tiny{[GM$|$GM]}}
                \texttt{yes} \\
            }
        }
    }
    & & \\ \\

    \theutterance \stepcounter{utterance}  
    & & & \multicolumn{2}{p{0.3\linewidth}}{
        \cellcolor[rgb]{0.9,0.9,0.9}{
            \makecell[{{p{\linewidth}}}]{
                \texttt{\tiny{[GM$|$GM]}}
                \texttt{game\_result = LOSE} \\
            }
        }
    }
    & & \\ \\

\end{supertabular}
}

\end{document}
